\newpage
\section{误差分析}

\subsection{预测模型误差}
\begin{enumerate}
    \item[(1)] \textbf{图结构构建}：邻接矩阵的构建方式（距离阈值vs K近邻）影响预测精度，我们实验发现混合策略效果最佳。
    \item[(2)] \textbf{输入特征}：仅使用历史流量数据，未纳入天气、事件等外生变量。
\end{enumerate}
MAE在正常交通条件下为12.3辆/5min，在拥堵条件下上升至28.7辆/5min。

\subsection{控制策略误差}
\begin{enumerate}
    \item[(1)] \textbf{协同假设}：分布式控制假设相邻路口信息完全可达，实际存在通信延迟和丢包。
    \item[(2)] \textbf{模型失配}：信号控制模型简化了排队演化过程，高峰期的溢出效应未被充分刻画。
\end{enumerate}
在10\%数据扰动下，控制策略的鲁棒性验证通过，总延误增加不超过8\%。
